% !TEX root = ../main.tex
\section{Decisiones solicitadas}
\label{sec:decision}
La fase preparatoria requiere definir cinco aspectos antes de seleccionar proveedor y concretar el piloto:
\begin{enumerate}
\item \textbf{Política de datos y nivel de aseguramiento.} Determinar la arquitectura admisible, el alcance ENS y la documentación exigida. El alojamiento local también requiere seguridad, mantenimiento y evaluación clínica.
\item \textbf{Servicios y participantes.} Identificar áreas prioritarias, tipos de consulta, profesionales y responsables clínico y técnico. Esta selección orientará los requisitos lingüísticos y de captura.
\item \textbf{Preparación del conjunto de evaluación.} Resolver si procede grabar consultas reales. Protección de Datos y el comité competente, cuando corresponda, deberán definir finalidad, base jurídica, información, consentimiento exigible, acceso, conservación y reutilización~\cite{rgpd2016}.
\item \textbf{Presupuesto y alcance comercial.} Fijar un límite para proveedores, equipos, integración, soporte y evaluación clínica. Las tarifas individuales no equivalen a un presupuesto hospitalario (apéndice~\ref{app:precios}).
\item \textbf{Infraestructura y trazabilidad.} Precisar equipos, sistemas, navegadores, periféricos y calidad de red. Decidir si basta con audio y nota o se requieren transcripción, versiones, correcciones y registros de transferencia y error, con finalidad y conservación definidas.
\end{enumerate}

Estas definiciones permitirán solicitar propuestas comparables y elaborar el protocolo. El siguiente hito será decidir si procede el piloto y en qué condiciones; la preparación no implica adjudicación ni autorización de uso asistencial.
